%----------------------------------------------------------------------------------------
%	PACKAGES AND OTHER DOCUMENT CONFIGURATIONS
%----------------------------------------------------------------------------------------

\documentclass[
    10pt,
]{style} % Use the style class

\usepackage{ebgaramond}
\usepackage{hyperref} % Use the EB Garamond font

%------------------------------------------------

\name{Hugo Pereira}

\address{+33626294012 \\ \texttt{hugo0807pereira@gmail.com}} % Contact information

%----------------------------------------------------------------------------------------

\begin{document}

    \begin{center}
        \url{https://www.linkedin.com/in/hugopereira75/}
    \end{center}

%----------------------------------------------------------------------------------------
%	WORK EXPERIENCE SECTION
%----------------------------------------------------------------------------------------

    \begin{rSection}{Expérience professionnelle}

        \begin{rSubsection}{Airbus Defence \& Space}{Février 2026 - Mai 2026}{Stage temps plein – Ingénieur cybersécurité (programme A400M)}{Getafe, Espagne}
            \item Stage de 3 mois en cybersécurité sur le programme d'avion de transport militaire A400M.
            \item Immersion dans l’ingénierie de systèmes sécurisés au sein d’un environnement défense fortement réglementé et international.
        \end{rSubsection}

        \begin{rSubsection}{Airbus Defence \& Space}{Août 2023 - Présent}{Alternance – Ingénieur logiciel}{Paris, France}
            \item Amélioration et évolution d’un plugin d’ingénierie système (MBSE) sur le logiciel CAMEO, développé en Java et utilisé par de nombreux ingénieurs d’Airbus.
            \item Mise en place et pilotage de la chaîne d’Intégration/Déploiement Continus (CI/CD) du projet, avec des pratiques DevOps : Jenkins, Docker, Git, Jira \& Confluence.
            \item Optimisation du système MBSE pour l’aligner précisément sur les besoins des architectes et faciliter son adoption.
        \end{rSubsection}

%------------------------------------------------

        \begin{rSubsection}{Kardham Digital}{Septembre 2022 - Août 2023}{Alternance – Développeur, Assistant Chef de Projet Digital}{Paris, France}
            \item Développement de solutions full-stack en Python, HTML, CSS, JavaScript, C\#, Java et Django REST API.
            \item Réalisation de projets pour le pôle Recherche \& Développement (prévision de présence, développement IoT, optimisation d’espaces, analyse de données bâtimentaires).
            \item Analyse de données avec SQL, Pandas, Streamlit et Elasticsearch afin de produire des statistiques de vie du bâtiment et des graphiques personnalisés pour les clients.
        \end{rSubsection}

%------------------------------------------------

        \begin{rSubsection}{Kardham Digital}{Juin 2022 - Août 2022}{Stage – Assistant Chef de Projet Digital}{Paris, France}
            \item Participation au suivi (BUILD \& RUN) de projets informatiques (gestion de projet avec SCRUM).
            \item Déplacements sur sites clients partout en France pour organiser le déploiement d’infrastructures IT : installation de routeurs, pare-feux, switches, capteurs et éclairages connectés.
            \item Administration de systèmes et hébergement de sites/projets web via Apache2 sur serveurs Debian.
        \end{rSubsection}

%------------------------------------------------

        \begin{rSubsection}{Université Paris Cité}{Octobre 2022 - Présent}{CDD – Tuteur étudiant}{Paris, France}
            \item Cours hebdomadaires et tutorat pour la section handicap de l’université.
        \end{rSubsection}

    \end{rSection}

%----------------------------------------------------------------------------------------
%	EDUCATION SECTION
%----------------------------------------------------------------------------------------

    \begin{rSection}{Formation}

        \textbf{Université de Technologie de Compiègne (UTC) - Sorbonne} \hfill \textit{Juin 2026} \\
        Diplôme d’ingénieur en génie logiciel \\
        Moyenne générale : 5/5 \\
        Spécialisation en POO, conception logicielle, algorithmique, structures de données, recherche opérationnelle \& optimisation combinatoire

        \textbf{Université Paris Cité} \hfill \textit{Juin 2023} \\
        DUT Informatique (diplôme universitaire de technologie) \\
        Major de promotion

    \end{rSection}

%----------------------------------------------------------------------------------------
%	TECHNICAL STRENGTHS SECTION
%----------------------------------------------------------------------------------------

    \begin{rSection}{Compétences techniques}

        \begin{tabular}{@{} >{\bfseries}l @{\hspace{6ex}} l @{}}
            Langages de programmation & Java, Swift, Objective-C, C\#, C, C++, Python, JS, HTML, CSS, VB .NET \\
            Systèmes & GNU/Linux (Debian, Ubuntu), environnement serveur, scripting shell, Ansible \\
            Cybersécurité & VPN, durcissement systèmes, gestion des accès, cyptographie, contexte défense (A400M) \\
            Outils & Vim, Linux, Git, Docker, Jenkins, Jira, Confluence, Trello, Slack \\
        \end{tabular}

    \end{rSection}

\end{document}